\documentclass[11pt]{article}
\usepackage[margin=.82in]{geometry}
\usepackage{amsmath,amssymb,booktabs,hyperref}
\hypersetup{colorlinks=true,urlcolor=blue}
\newcommand{\F}{\mathbb F}
\newcommand{\Q}{\mathbb Q}
\newcommand{\Z}{\mathbb Z}
\title{\vspace{-1.5em}An exact rank rule for sparse roots\\ in closed \(p\)-adic disks}
\author{ProveIt GRH article supplement}
\date{27 November 2026}
\begin{document}
\maketitle

A finite rank calculation determines the maximum number of roots
for a specified set of at most \(p\) exponents, at every depth where
the \(p\)-adic exponential converges. The answer can distinguish
supports with the same residue multiplicities. At \(p=7\) and
depth \(1/5\), the supports \((0,1,2,7,8,14)\) and \((0,1,3,7,8,14)\)
give exact maxima \(14\) and \(8\), respectively.
The same rule yields the exact unrestricted seven-term profile
at \(p=7\), including the intermediate maximum \(9\).

\section{The rank rule}
Let \(E\) consist of \(2\le t\le p\) distinct nonnegative integers.
Allow finite extensions of \(\Q_p\) for the coefficients and for the
roots, normalize \(v_p(p)=1\), and fix a nonzero center \(A\) in a
finite extension. Count roots in \(v_p(X/A-1)\ge r>1/(p-1)\).
Put
\[
 R=\F_p[T]/\prod_{n\in E}(T-\bar n),\qquad
 \mathcal K=\{i+pj:i,j\ge0,\ i+j\le t-2\}\cup\{t-1\},
 \qquad w_r(k)=rk-\lfloor k/p\rfloor .
\]
The maximum is the largest of \(t-1\) and those \(k\in\mathcal K\)
with \(w_r(k)\le w_r(t-1)\) for which
\[
 T^k\notin \operatorname{span}_{\F_p}\{T^\ell:
       \ell\in\mathcal K,\ w_r(\ell)<w_r(k)\}.
 \tag{1}
\]
The strict inequality is essential: tied weights must be tested
against the same lower-weight span. There are at most
\(1+t(t-1)/2\) vectors. The rule is a rational staircase, computed
by sorting pairwise weight crossings. The upper bound counts
multiplicity, and the maximum is attained by distinct roots,
with all terms present. At center \(1\) the coefficients may be
chosen rational. After dividing out the smallest monomial, the
witness may even be chosen irreducible over \(\Q\).

\paragraph{Necessity.}
Start with rational \(r\). The
\href{https://github.com/VladimirReshetnikov/ProveIt/blob/dd990b27a22290c756c1bd4adce1d971bf7e7821/Analysis/GeneralizedRiemannHypothesis/Article/sections/smd.tex#L57}{ProveIt local transformation}
\(X=A\exp(bY)\), \(v_p(b)=r\), gives a restricted power series
with \(k\)-th coefficient \(M_k b^k/k!\), where
\(M_k=\sum_{n\in E}c_n A^n n^k\). After dividing by a moment of
least valuation among the first \(t\), all moments are integral by
the recurrence with polynomial \(\prod_{n\in E}(T-n)\);
at least one early moment is a unit by Vandermonde.
For general \(k\), write \(w(k)=rk-v_p(k!)\).
The number of roots with multiplicity is the largest minimizing
coefficient index, by Weierstrass preparation on the closed unit disk.
\newpage
A minimizing \(k\ge t\) must have \(w(k)\le w(t-1)\).
Legendre's formula gives
\[
 w(k)=(r-1/(p-1))k+S_p(k)/(p-1)
\]
with digit sum \(S_p(k)\). Thus \(S_p(k)\le t-2\), \(k\ge p\),
and \(w(k)\ge w(p)\). Also \(w(t-1)-w(p)=1-(p-t+1)r<1\).
The normalized moment at such a \(k\) has valuation below one.

On each factor of \(R\) at residue \(a\), \(U=T^p=a\),
because the multiplicity is at most \(p\), and \(T^{p^j}=U\)
for \(j\ge1\). Compressing higher digits to the \(p\)-digit keeps
the same vector and digit sum, and strictly lowers the weight
unless the index already belongs to \(\mathcal K\).
On \(\mathcal K\), \(w=w_r\) because \(k<p^2\).
If the resulting vector is spanned by strictly lower weights,
lift the relation modulo \(p\) to the integral quotient.
Since the moment at \(k\) has valuation below one, some
contributing lower-weight moment has no larger valuation.
Its series coefficient has strictly smaller valuation,
a contradiction. This also explains why ramified coefficient
fields do not alter the count.

\paragraph{Sufficiency and witnesses.}
For a passing \(k>t-1\), choose a residue functional taking value
\(1\) on \(T^k\) and annihilating the lower-weight span.
Lift it to an integral functional \(L\) in the power basis of
the actual exponent quotient. At basis indices of weight below
\(w(k)\), prescribe exact zero, which lifts their residue values.
Any remaining lower-weight index is at least \(p\).
Its digit sum is at most \(t-2\), so compression puts its vector
in the lower-weight span, giving a moment divisible by \(p\).
Its coefficient therefore has valuation at least \(1+w(p)>w(k)\).
All other moments are integral, so the \(k\)-th coefficient
attains the minimum. At center \(1\), recover rational coefficients
by applying \(L\) to the Lagrange cardinal polynomials on \(E\).
For the default \(t-1\), use the functional that extracts the
\(T^{t-1}\)-coefficient in the power basis, guaranteeing order
at least \(t-1\) at \(1\). The upper bound gives equality when
no larger index passes.

To obtain distinct roots and exact support, perturb the rational
coefficients by sufficiently high \(p\)-powers. This preserves the
last minimizing coefficient index, since each perturbation
coefficient in the series has valuation at least the precision
plus \(w(j)\ge0\). Clear denominators, first ensuring every support
coefficient is nonzero. Having fixed the leading coefficient,
use CRT to impose Eisenstein at a different prime not dividing
it, after factoring out \(X^{\min E}\). This preserves support
and makes the nonzero roots simple. For general \(A\), scale
the variable; rational irreducibility is asserted at \(A=1\).

At a crossing, all strictly lower inequalities there already
hold just to the left, while candidate eligibility has a closed
inequality. Thus the left-hand maximum remains feasible at
equality. Monotonicity gives the opposite inequality. For an
irrational depth the adjacent rational samples in the same
cell give matching upper and lower bounds.
\newpage
\section{Two supports and the seven-term profile at \(p=7\)}
The following tables use \(v_7(X-1)\ge r>1/6\).
All thresholds retain the larger count at equality.
\begin{center}
\begin{tabular}{ll}
\toprule
Support & Successive maxima and transition depths\\
\midrule
\((0,1,2,7,8,14)\) & \(14\xrightarrow{\,1/5\,}7\xrightarrow{\,1/2\,}5\)\\
\((0,1,3,7,8,14)\) & \(14\xrightarrow{\,2/11\,}8\xrightarrow{\,1/4\,}7
          \xrightarrow{\,1/2\,}5\)\\
\bottomrule
\end{tabular}
\end{center}
For an explicit audit at \(r=1/5\), the corresponding integer
coefficient vectors (in the displayed exponent orders) are
\[
 (-17792,50400,-36946,14976,-10465,170),\qquad
 (86,504,-343,255,-161,2).
\]
Direct binomial expansion at \(X=1\) confirms their counts.

Maximizing over \emph{all} supports of at most seven terms at \(p=7\),
the exact profile is
\[
       14\xrightarrow{\,1/5\,}9\xrightarrow{\,1/3\,}8
          \xrightarrow{\,1/2\,}7\xrightarrow{\,1\,}6 .
\]
This is a finite classification using (1), not an experimental
extrapolation. Use seven slots (pad with zero coefficients) and
enumerate all \(\binom{13}{7}=1716\) residue multisets. Any such
multiset is realized by distinct nonnegative integer exponents.
The following supports realize the larger count at the indicated
endpoints:
\begin{center}
\begin{tabular}{rrl}
\toprule
Depth & Count & Support\\
\midrule
\(1/5\)&14&\((0,1,2,7,8,14,21)\)\\
\(1/3\)&9&\((0,1,2,4,7,14,21)\)\\
\(1/2\)&8&\((0,1,6,7,14,21,28)\)\\
\(1\)&7&\((0,1,7,14,21,28,35)\)\\
\bottomrule
\end{tabular}
\end{center}

\paragraph{Replay and scope.}
The script computes each residue vector by its truncated
expansions at the individual residues and does exact elimination.
It speeds up the exhaustive seven-slot classification by taking
one representative of each affine residue orbit. Under
\(T\mapsto aT+b,\ U\mapsto aU+b\), where \(a\ne0\), the
expansion of \(T^iU^j\) changes only by a nonzero scalar and
terms of strictly smaller weight: reducing \(i\) costs \(r>0\),
and reducing \(j\) costs \(pr-1>0\). The same holds for the
extra power \(T^{t-1}\). This preserves the rank test.
Endpoint witnesses are reconstructed by exact rational
arithmetic and checked independently by Newton weights in \(X-1\).
The replay also checks 7,218 seeded support/depth cases.

This extends the article's local estimate in the
prescribed-support setting. For the classical background on
general exponential polynomials see
\href{https://doi.org/10.1112/jlms/s2-36.1.1}{van der Poorten--Rumely}.
No RH or GRH hypothesis is used. The theorem is a written proof
with exact replays, not a Lean formalization or priority claim.
\end{document}
